\documentclass[../cours_1re.tex]{subfiles}
\begin{document}
\setcounter{chapter}{25}
\chapter{Complément : construire exponentielle et logarithme}
\label{ch:26}
\begin{objectif}
Complément : démontrer l'existence de l'exponentielle, puis construire son inverse logarithme.
Prérequis : chapitre 25, sommes, factorielles. Ce chapitre ferme l'admission explicite du chapitre 12.
\end{objectif}
\section{Approcher par des polynômes}
Posons $P_N(x)=\sum_{k=0}^N x^k/k!$.
Leurs dérivées satisfont $P_N'=P_{N-1}$ pour $N\ge1$ et $P_N(0)=1$.
Fixons $M>0$. Dès que $k+1\ge2M$,
\[
\frac{M^{k+1}/(k+1)!}{M^k/k!}=\frac M{k+1}\le\frac12.
\]
Les termes suivants sont donc dominés par une suite géométrique.
Pour $N+2\ge2M$ et $|x|\le M$, la somme des termes d'indice au moins $N+1$
est majorée en valeur absolue par
$2M^{N+1}/(N+1)!$, qui tend vers $0$.
Cela rend $(P_N(x))$ de Cauchy pour chaque $x$ et donne une limite
\[
E(x)=\sum_{k=0}^{\infty}\frac{x^k}{k!},
\qquad
\sup_{|x|\le M}|E(x)-P_N(x)|\le\frac{2M^{N+1}}{(N+1)!}.
\]
La convergence est \emph{uniforme} sur chaque segment $[-M,M]$ :
un même rang contrôle toutes ses abscisses.

\section{Pourquoi on peut passer à la limite}
Une limite uniforme de fonctions continues est continue.
En effet, pour un point $a$ et une tolérance $\varepsilon>0$, choisir $N$
tel que $|E-P_N|<\varepsilon/3$ sur un voisinage fermé de $a$.
La continuité de $P_N$ donne $|P_N(x)-P_N(a)|<\varepsilon/3$ pour $x$
assez proche. Ajouter les trois différences donne $|E(x)-E(a)|<\varepsilon$.

L'intégration des polynômes donne
$P_N(x)=1+\int_0^xP_{N-1}(t)\,dt$.
La borne
\[
\left|\int_0^x(P_{N-1}(t)-E(t))\,dt\right|
\le |x|\sup_{|t|\le |x|}|P_{N-1}(t)-E(t)|
\]
permet de passer à la limite. Donc $E(x)=1+\int_0^xE(t)\,dt$.
Comme $E$ est continue, le théorème fondamental du chapitre 25 donne
$E'=E$ et $E(0)=1$.
L'existence de l'exponentielle est maintenant démontrée. L'unicité,
la positivité et la loi d'addition ont été prouvées au chapitre 12 à partir
de cette seule équation différentielle.

\begin{attention}
On n'a pas simplement « dérivé une somme infinie ».
La convergence uniforme et l'intégration ont justifié le passage à la limite.
Une convergence seulement point par point ne suffit pas à préserver toutes
les propriétés : $x^n$ sur $[0,1]$ converge vers une fonction discontinue.
\end{attention}

\section{Calculer le nombre \texorpdfstring{$e$}{e}}
$e=E(1)=\sum_{k=0}^\infty1/k!$.
Pour $N\ge1$,
$0<e-\sum_{k=0}^N1/k!\le2/(N+1)!$.
Par exemple, $N=9$ donne une erreur inférieure à
$2/10!\simeq5{,}52\times10^{-7}$.
La formule résulte de la même domination géométrique, et fournit
une précision garantie en arithmétique exacte.

\section{Définir le logarithme par une bijection}
L'exponentielle est continue, strictement croissante, de limite $0$
en $-\infty$ et $+\infty$ en $+\infty$ (chapitre 12).
Le théorème des valeurs intermédiaires assure que chaque $y>0$
a un unique antécédent. On le note $\ln y$ :
\[
\exp(\ln y)=y\quad(y>0),\qquad \ln(\exp x)=x\quad(x\in\R).
\]
L'inverse est strictement croissante. Elle est continue :
si $y_0=e^{x_0}$ et $\varepsilon>0$, les nombres
$e^{x_0-\varepsilon}<y_0<e^{x_0+\varepsilon}$ encadrent un voisinage
de $y_0$ ; leurs images par $\ln$ encadrent $\ln y$ entre
$x_0-\varepsilon$ et $x_0+\varepsilon$.

\begin{demonstration}[dérivée du logarithme]
Pour $y+h>0$, posons $x=\ln y$ et $k=\ln(y+h)-\ln y$.
Si $h\ne0$, alors $k\ne0$ ; par continuité $k\to0$.
Comme $h=e^{x+k}-e^x$,
\[
\frac{\ln(y+h)-\ln y}{h}
=\frac1{(e^{x+k}-e^x)/k}\longrightarrow\frac1{e^x}=\frac1y.
\]
\end{demonstration}
L'injectivité de l'exponentielle et sa loi d'addition donnent
$\ln(ab)=\ln a+\ln b$, $\ln(1/a)=-\ln a$ pour $a,b>0$.
On a $\ln1=0$, $\ln e=1$ et $\ln(a^r)=r\ln a$ en définissant
$a^r=\exp(r\ln a)$ pour un réel $r$.
Ce prolongement est cohérent avec les puissances entières et rationnelles.

\section{Deux encadrements utiles}
Pour $t>-1$, $\ln(1+t)\le t$ : la fonction
$t-\ln(1+t)$ a pour dérivée $t/(1+t)$ et un minimum nul en $0$.
Pour $t\ge0$, on a aussi $\ln(1+t)\ge t/(1+t)$ :
la différence a pour dérivée $t/(1+t)^2\ge0$ et vaut $0$ en $0$.
Avec $t=1/n$ :
\[
\frac{n}{n+1}\le n\ln(1+1/n)\le1.
\]
Les gendarmes puis la continuité de l'exponentielle prouvent
$(1+1/n)^n\to e$. Ce n'est donc pas une définition concurrente non justifiée.

Enfin $\int_1^xdt/t=\ln x$, car les deux fonctions ont la même dérivée
et s'annulent en $1$. Le logarithme mesure ainsi une aire et transforme
les produits en sommes : deux descriptions du même objet.

\begin{finchapitre}
\section{Utilité et prolongements}
Les séries donnent des valeurs approchées avec contrôle d'erreur ;
les fonctions inverses résolvent des équations, par exemple
$e^{-kt}=1/2\iff t=\ln2/k$ pour $k>0$.

La convergence uniforme sera centrale pour les séries de fonctions en licence.
Les équations intégrales et les opérateurs de limite deviendront des outils
de construction de solutions en analyse au niveau master.
\begin{retenir}
La série de l'exponentielle converge avec un contrôle uniforme sur chaque segment.
L'intégration justifie sa dérivabilité. Le logarithme est l'inverse de l'exponentielle sur les réels positifs.
\end{retenir}
\end{finchapitre}
\end{document}
